\begin{frame}{Peer-Review 결실: 예시 2와 3점의 형태}
  \begin{alertblock}{예시 2 — ``\textbf{왜 딥러닝인가}'' 질문
    (16.2 추가 항목)}
    \scriptsize 발표: ``고객 이탈 예측(\textbf{정형 테이블
    데이터})에 \textbf{Transformer} 적용, test 81\%. 최신
    아키텍처라 기대가 컸다.''
  \end{alertblock}
  \vskip0.2em
  \scriptsize
  \begin{tabular}{@{}p{2.6cm}p{12.1cm}@{}}
    \toprule
    단계 & 리뷰 \\
    \midrule
    1점(감상평) & ``최신 모델을 적용하셨네요. 81\%는 좋은
      결과입니다.'' \\
    2점(일반 지적) & ``정형 데이터에 Transformer를 쓴 근거가
      \textbf{약해} 보입니다.'' \\
    3점(반박 가능) & ``정형 데이터에서 \textbf{GBDT(Ch07.3)}는
      신경망과 \textbf{동일하거나 높은} 성능이 흔합니다 — 같은
      데이터의 GBDT \textbf{val F1은 얼마인가요}(관찰 요구)?
      Transformer가 여기서 주는 \textbf{구체적 구조적 장점}(규모,
      상호작용 패턴)이 \textbf{무엇이고}, 그 장점이 \textbf{숫자
      어디에} 보입니까(개념)? \textbf{새 계산 없이} 답할 수
      없다면 ``딥러닝이 필요했다''는 결론이 성립하지 않습니다.'' \\
    \bottomrule
  \end{tabular}
  \vskip0.4em
  \begin{block}{두 예시의 3점 질문은 \textbf{같은 형태}}
    \textbf{관찰}(발표에서 본 것) + \textbf{개념}(장 번호) +
    \textbf{왜 문제인지} + \textbf{구체적 산출물 요구} — 네 칸
    이 \textbf{모두} 채워져야 ``좋은 지적입니다''로 넘길 수 없는
    질문(8.2 FAQ).
  \end{block}
  \vskip0.3em
  {\footnotesize 이 학기 동안: 발표자 2회(8.1, 16.1), 리뷰어
  2회(8.2, 16.2). 8.2의 \textbf{``반증 가능성(falsifiability)''}
  이 이제 낯선 단어가 아니라면 — \textbf{이 절차는 끝난 것이다}.}
\end{frame}
